
\medskip
\noindent \SGARefTarget{sga7.sga7i.eviii.n2-3-19}{2.3.19}.  \underline{Another proof of \SGARefLink{sga7.sga7i.eviii.n2-3-11}{2.3.11}} (by
Deligne).

\medskip
\noindent a) Reduction to $P=\mathbb Z/n\mathbb Z$ and to
$T=(\mathrm{Sets})$.

\medskip
\noindent b) Let $\partial$ be the following arrow from $nQ$ to
$\,{}_nQ[1]$:
\[
 nQ\xleftarrow[\text{``$n$''}]{\ \sim\ }
 [\,{}_nQ\longrightarrow Q\,]
 \longrightarrow{}_nQ[1].
\]
The following diagram is anticommutative:
\[
\begin{tikzcd}[row sep=huge,column sep=huge]
 nQ[0]\arrow[r,"1"]\arrow[d,"\partial"']
   &{[Q\xrightarrow{n}Q]}\\
 {}_nQ[1]\arrow[ru,"1"']&{.}
\end{tikzcd}
\]
Indeed, the arrows
\[
\begin{tikzcd}[column sep=small,row sep=large]
 {[{}_nQ}\arrow[r]\arrow[dd,"0"']&{Q]}\arrow[dd,"n"]
   &&{[{}_nQ}\arrow[r]\arrow[dd,"1"']&{Q]}\arrow[dd,"0"']
   &\text{are opposite}
   &{[{}_nQ}\arrow[r]&{Q]}\arrow[ddl,"1"]\\
 &&\text{and}&&&\text{by means of}\quad H\,:&&\\
 {[Q}\arrow[r,"n"]&{Q]}&&{[Q}\arrow[r,"n"]&{Q]}
   &&{[Q}\arrow[r,"n"]&{Q]}
\end{tikzcd}
\]

\medskip
\noindent c) Diagram \SGARefLink{sga7.sga7i.eviii.n2-3-10}{2.3.10} is
\[
\begin{tikzcd}[column sep=small,row sep=small]
 &&&H^1\mathrm{RHom}
   (\mathbb Z\mathbin{\overset{\mathbb L}{\otimes}}Q,G)
   \arrow[ddd]\\
 &&&+\\
 H^1\mathrm{RHom}
   (\mathbb Z/n\mathbb Z\mathbin{\overset{\mathbb L}{\otimes}}Q,G)
   \arrow[rr]\arrow[dd]\arrow[uurrr]
 &&H^1\mathrm{RHom}
   (\mathbb Z/n\mathbb Z,\mathrm{RHom}(Q,G))
   \arrow[dd]&\\
 &&&\mathrm{Ext}^1(Q,G)\arrow[ddl]\\
 \mathrm{Hom}(\mathrm{Tor}_1(\mathbb Z/n\mathbb Z,Q),G)\arrow[d]
 &|[draw,circle,inner sep=1pt]|{-\,?}
 &\mathrm{Hom}(\mathbb Z/n\mathbb Z,\mathrm{Ext}^1(Q,G))
   \arrow[d]&\\
 \mathrm{Hom}({}_nQ,G)\arrow[rr]
 &&\mathrm{Ext}^1(nQ,G)&
\end{tikzcd}
\]
and:

\smallskip
\noindent c1) We have an isomorphism
\[
 \mathbb Z/n\mathbb Z
 \mathbin{\overset{\mathbb L}{\otimes}}Q
 \simeq[Q\xrightarrow{n}Q],
\]
and, under this isomorphism, the first composite vertical arrow is
induced by
\[
 {}_nQ[1]\longrightarrow[Q\xrightarrow{n}Q].
\]

\smallskip
\noindent c2) The second horizontal arrow is the transpose of the
arrow $\partial$ displayed above in b).

\smallskip
\noindent c3) If the diagram is traversed clockwise, it is clear that
one obtains the arrow induced by
\[
 nQ[0]\longrightarrow[Q\xrightarrow{n}Q].
\]
This proves \SGARefLink{sga7.sga7i.eviii.n2-3-11}{2.3.11}.
